\mxunit{automatismes-ch07}{Proportions et évolutions}

\begin{fluencysheet}{mar.~2~févr.~2027}{Proportions et pourcentages\,; proportion de proportion}
\mxfq{Échauffement · Flash}{PRO}{Calculer $\pc{15}$ de $240$.}{$\num{0,15}\times240=36$.}{}
\mxfq{Cœur · Méthode}{PRO}{$\pc{30}$ des élèves sont internes et $\pc{40}$ des internes pratiquent la danse. Quelle proportion des élèves sont des internes qui dansent\,?}{$\num{0,3}\times\num{0,4}=\num{0,12}=\pc{12}$~: proportion de proportion (on multiplie).}{}
\mxfq{Cœur + · Problème}{PRO}{Un lycée compte $800$ élèves dont $\pc{55}$ de filles\,; $\pc{20}$ des filles font du latin. Combien y a-t-il de filles latinistes\,?}{$800\times\num{0,55}=440$ filles\,; $440\times\num{0,20}=88$.}{}
\mxfq{Retour · Flash}{QP}{$3$ cahiers coûtent $\num{4,50}$ €. Quel est le prix de $7$ cahiers\,?}{$\num{1,50}\times7=\num{10,50}$ €.}{}
\mxfq{Fil rouge · Vrai/faux}{EQ}{Vrai ou faux\,? Justifier. Le nombre $-2$ est solution de l'équation $3x+1=-5$.}{Vrai. $3\times(-2)+1=-6+1=-5$.}{}
\end{fluencysheet}

\begin{fluencysheet}{jeu.~4~févr.~2027}{Variation absolue, coefficient multiplicateur, taux d'évolution}
\mxfq{Échauffement · Flash}{PRO}{Un prix passe de $80$ Ar à $92$ Ar. Quelle est la variation absolue\,?}{$92-80=12$ Ar.}{}
\mxfq{Cœur · Méthode}{PRO}{Pour l'évolution de $80$ Ar à $92$ Ar, donner le taux d'évolution et le coefficient multiplicateur.}{$t=\dfrac{12}{80}=\num{0,15}=\pc{15}$\,; $CM=1+t=\num{1,15}$.}{}
\mxfq{Cœur + · QCM}{PRO}{Une baisse de $\pc{20}$ correspond au coefficient multiplicateur~:\mxchoices{$\num{0,2}$}{$-\num{0,8}$}{$\num{1,2}$}{$\num{0,8}$}}{D.\ $\num{0,8}$}{$CM=1+t=1-\num{0,20}$.}
\mxfq{Retour · Flash}{STA}{Les températures relevées (en °C) sont $-2,\,0,\,3,\,5,\,9$. Quelles sont leur moyenne et leur étendue\,?}{Moyenne $3$ °C\,; étendue $9-(-2)=11$ °C.}{}
\mxfq{Fil rouge · Flash}{SOL}{Calculer le volume et l'aire totale d'un cube d'arête $3$ cm.}{$V=3^3=27$ cm$^3$\,; $A=6\times3^2=54$ cm$^2$.}{}
\end{fluencysheet}

\begin{fluencysheet}{lun.~8~févr.~2027}{Évolutions successives}
\mxfq{Échauffement · Flash}{CAL}{Calculer $\num{1,2}\times\num{0,9}$.}{$\num{1,08}$.}{}
\mxfq{Cœur · Méthode}{PRO}{Une hausse de $\pc{20}$ est suivie d'une baisse de $\pc{10}$. Quel est le taux global\,?}{$\num{1,2}\times\num{0,9}=\num{1,08}$~: hausse globale de $\pc{8}$.}{Les taux ne s'additionnent pas\,; ce sont les coefficients qui se multiplient.}
\mxfq{Cœur + · Vrai/faux}{PRO}{Vrai ou faux\,? Justifier. Une hausse de $\pc{50}$ suivie d'une baisse de $\pc{50}$ ramène à la valeur initiale.}{Faux. $\num{1,5}\times\num{0,5}=\num{0,75}$~: baisse globale de $\pc{25}$.}{}
\mxfq{Retour · Flash}{SUB}{Calculer $x^2-x$ pour $x=\dfrac12$.}{$-\dfrac{1}{4}$.}{}
\mxfq{Fil rouge · Rédiger}{DEM}{$ABCD$ est un parallélogramme tel que $AC=BD$. Quelle est sa nature\,? Justifier en une phrase.}{Rectangle~: un parallélogramme dont les diagonales ont même longueur est un rectangle.}{}
\end{fluencysheet}

\begin{fluencysheet}{mar.~9~févr.~2027}{Évolution réciproque\,; proportion et évolution}
\mxfq{Échauffement · Flash}{CAL}{Calculer $\dfrac{1}{\num{1,25}}$.}{$\num{0,8}$.}{}
\mxfq{Cœur · Méthode}{PRO}{Après une hausse de $\pc{25}$, quelle baisse ramène à la valeur initiale\,?}{$CM_{\text{réc}}=\dfrac{1}{\num{1,25}}=\num{0,8}$~: baisse de $\pc{20}$ (et non $\pc{25}$).}{Piège~: croire qu'une hausse de $\pc{25}$ se compense par une baisse de $\pc{25}$.}
\mxfq{Cœur + · Interpréter}{PRO}{Sur $250$ candidats, $150$ sont admis\,; l'an dernier il y avait $200$ candidats. Calculer la proportion d'admis puis le taux d'évolution du nombre de candidats.}{Proportion $\dfrac{150}{250}=\pc{60}$ (partie sur tout)\,; évolution $\dfrac{250-200}{200}=\pc{25}$ (variation relative).}{}
\mxfq{Retour · Tableau}{IMG}{Voici un tableau de valeurs de $f$~:\par\smallskip\mxarraystretch{1.15}\begin{tabular}{|c|c|c|c|c|}\hline
$x$ & 0 & 1 & 2 & 3 \\ \hline
$f(x)$ & 4 & 1 & 0 & 1 \\ \hline
\end{tabular}\par Quelle est l'image de $2$\,? Quels sont les antécédents de $1$\,?}{$f(2)=0$\,; antécédents de $1$~: $1$ et $3$.}{}
\mxfq{Fil rouge · QCM}{REP}{Sur le globe terrestre, deux points d'un même parallèle ont~:\par\smallskip A.~la même latitude\quad B.~la même longitude\par C.~la même distance à Greenwich\quad D.~la même altitude}{A.\ la même latitude}{Latitude~: angle nord/sud par rapport à l'équateur\,; longitude~: angle est/ouest par rapport au méridien de Greenwich.}
\end{fluencysheet}
